\chapter{Métodos e recursos}

Neste capítulo são descritos os procedimentos metodológicos adotados para o alcance dos objetivos propostos, abordando a abordagem de pesquisa escolhida, os instrumentos utilizados na coleta e análise de informações e as ferramentas empregadas no desenvolvimento da API.

\section{Visão geral}

Para o desenvolvimento da pesquisa adotou-se a abordagem qualitativa, que se destaca por enfatizar a análise de dados não numéricos, como documentos, normativas e literatura especializada, buscando compreender o significado e a complexidade dos fenômenos investigados.

\begin{citacao}
A pesquisa qualitativa trabalha com o universo de significados, motivos, aspirações, crenças, valores e atitudes, o que corresponde a um espaço mais profundo das relações, dos processos e dos fenômenos que não pode ser reduzido à operacionalização de variáveis. \cite[p.~21]{minayo2001}
\end{citacao}

Para a condução do trabalho, foram traçadas etapas sequenciais, ilustradas na Figura~\ref{fig:processo-tcc} e na Figura~\ref{fig:processo-api}.

\begin{figure}[H]
  \centering
  \includegraphics[width=0.95\textwidth]{figura01-processo-tcc}
  \caption{Diagrama de processos do desenvolvimento do TCC}
  \fonte{Elaborado pelos próprios autores (2025).}
  \label{fig:processo-tcc}
\end{figure}

\begin{figure}[H]
  \centering
  \includegraphics[width=0.95\textwidth]{figura02-processo-api}
  \caption{Diagrama de processos do Desenvolvimento da API}
  \fonte{Elaborado pelos próprios autores (2025).}
  \label{fig:processo-api}
\end{figure}

Inicialmente, definiu-se o escopo da pesquisa, o problema, as justificativas e os objetivos que fundamentam o Capítulo~1. Em seguida, conduziu-se o levantamento bibliográfico sobre OO, SOLID e DDD, aprofundados no Capítulo~2.

Paralelamente, realizou-se o mapeamento do processo de um hemocentro por meio de pesquisa documental e bibliográfica, buscando compreender o fluxo de trabalho, as regras de negócio e as necessidades do domínio da doação de sangue. Elaborou-se também um protocolo de mapeamento sistemático da literatura, cujos resultados são discutidos no Capítulo~3.

Com base no conhecimento adquirido nessas etapas, realizou-se a análise do domínio sob a ótica do DDD e a definição dos requisitos funcionais e não funcionais. Na sequência, elaborou-se o projeto da solução (casos de uso e diagramas de classes) e desenvolveu-se a API aplicando SOLID, com vistas à escalabilidade, manutenibilidade e testabilidade.

Por fim, os resultados foram avaliados por meio de testes unitários e de integração, seguidos da elaboração das conclusões do trabalho.

\section{Instrumentos e ferramentas}

Para a coleta e análise dos dados, foram utilizados instrumentos compatíveis com os objetivos e a abordagem metodológica da pesquisa, orientados pela necessidade de obter informações confiáveis e contextualizadas sobre o cenário da doação de sangue e sobre as tecnologias empregadas no desenvolvimento da API.

\subsection{Pesquisa exploratória-descritiva}

A presente pesquisa caracteriza-se como exploratória-descritiva, abordagem que, segundo \citeonline{gil2008}, proporciona maior familiaridade com o problema investigado e, simultaneamente, descreve as características do fenômeno estudado. Segundo \citeonline{prodanov2013}, essa combinação metodológica é especialmente adequada para estudos que envolvem questões sociais e de saúde pública, permitindo uma análise contextualizada e orientada para a intervenção.

O contexto de análise adotado foi o das normativas e do cenário de atuação dos hemocentros, com atenção às demandas da região de Campos dos Goytacazes, local de moradia e estudo dos pesquisadores. Essa escolha favoreceu maior familiaridade com os aspectos sociais, culturais e institucionais que envolvem a temática, além de facilitar o acesso a dados públicos e normativas das unidades hospitalares e hemocentros locais.

Segundo \citeonline{pressman2016}, o esforço inicial de compreensão do domínio e das necessidades reais dos \textit{stakeholders} é indispensável para criar uma base sólida de requisitos e evitar o desenvolvimento de sistemas inadequados à realidade operacional. Nesse sentido, utilizou-se a análise documental e bibliográfica como técnica de coleta de dados, por permitir profundidade na obtenção de informações normativas e procedimentais a partir de manuais do Ministério da Saúde, portarias vigentes e literatura acadêmica.

Segundo \citeonline{marconi2003}, a pesquisa documental é composta por etapas fundamentais --- como a seleção, leitura crítica e categorização dos dados --- que garantem a sistematização e a confiabilidade das informações obtidas. \citeonline{sommerville2011} reforça que a análise de documentos atua como instrumento de captura e construção do conhecimento, sendo essencial para a engenharia de requisitos em sistemas críticos, como os voltados à saúde pública.

O levantamento realizado a partir de documentos oficiais e estudos de caso sobre triagem, captação e gestão em hemocentros teve como objetivo compreender os fatores que influenciam a decisão de doar sangue, identificar barreiras e motivações entre potenciais doadores descritos na literatura, e subsidiar os requisitos da API proposta.

\subsection{Desenvolvimento da API}

O desenvolvimento da API foi orientado por boas práticas de engenharia de \textit{software}, com foco na escalabilidade, na manutenibilidade e na clareza estrutural do código. Adotou-se DDD e SOLID como diretrizes arquiteturais (Seções~\ref{sec:ddd} e~\ref{sec:solid}). A camada de segurança contempla autenticação baseada em \textit{JSON Web Token} (JWT), padrão amplamente adotado em APIs RESTful para autenticação sem estado (\textit{stateless}).

\subsubsection{Linguagem de programação (Java)}

A linguagem escolhida para a implementação foi o Java, em sua versão 17 (\textit{Long Term Support} --- LTS). Java é uma linguagem orientada a objetos, fortemente tipada e amplamente documentada na literatura \cite{deitel2017}. A tipagem estática do Java favorece a implementação de \textit{Value Objects} e \textit{Entities}, garantindo que erros de domínio sejam capturados em tempo de compilação. Em contrapartida, uma crítica comum ao Java é a sua verbosidade e o consumo de memória da \textit{Java Virtual Machine} (JVM), se comparado a linguagens como Go ou Python, conforme discutido por \citeonline{bloch2018}. Ainda assim, a robustez da linguagem compensa essas restrições para o contexto corporativo.

\subsubsection{\textit{Framework} de desenvolvimento (Spring Boot)}

Para agilizar o desenvolvimento e a configuração da infraestrutura, utilizou-se o \textit{framework} Spring Boot. Segundo \citeonline{walls2019}, o Spring Boot aplica o conceito de ``convenção sobre configuração'', permitindo que o desenvolvedor foque na lógica de negócio enquanto o \textit{framework} gerencia a injeção de dependência e a conexão com o banco de dados. Por outro lado, a facilidade do Spring pode esconder a complexidade subjacente, dificultando o \textit{debug} (depuração) caso o desenvolvedor não entenda o ciclo de vida dos \textit{beans} gerenciados pelo \textit{container}.

\subsubsection{Ambiente de desenvolvimento e versionamento (Visual Studio Code e Git)}

Para a codificação, utilizou-se o editor de código Visual Studio Code (VS Code), devido ao seu vasto ecossistema de extensões para Java e suporte nativo a terminais leves. Para o controle de versão, adotou-se o Git, hospedando o repositório na plataforma GitHub. Segundo \citeonline{chacon2014}, o versionamento distribuído é essencial para manter o histórico de alterações seguro e permitir ramificações (\textit{branches}) para testes de novas funcionalidades sem comprometer o código estável (\textit{main}).